% GENERATED FILE. Edit canonical lessons, then run npm run generate:books.
% canonical-source-hash: fnv1a64:5d9c9bdbe3f6aeff
% canonical-lessons: GE-C241-umwelt, GE-C241-unterricht, GE-C241-unterschied, GE-C241-ursache, GE-C241-verhalten

\chapter{The environment, Lessons, A difference, A cause, Behaviour}
\label{ch:ge-the-environment-lessons-a-difference-a-cause-behaviour}
\hlchaptermodality{241}

\begin{chapteropening}
\textbf{By the end of this chapter:} \emph{I can say the environment, lessons, a difference, a cause and behaviour.}

Complete the last lesson of chapter 241: I can say the environment, lessons, a difference, a cause and behaviour.
\end{chapteropening}

\section[die Umwelt]{die Umwelt — the environment}

\label{lesson:GE-C241-umwelt}



\begin{quote}
Before the new one: say the German for a surprise, then the German for the surroundings.
\end{quote}

\subsection*{You'll want to know: die Umwelt}
\textbf{die Umwelt} — \textquotedblleft{}the environment\textquotedblright{}.

\begin{cousinweb}
\textbf{Umwelt} is \emph{um-}, \textquotedblleft{}around\textquotedblright{}, + \emph{Welt}, \textquotedblleft{}world\textquotedblright{}: the world around us. \emph{Welt} was first \textquotedblleft{}the age of man\textquotedblright{}. (EtymDB 2.1)
\end{cousinweb}

\begin{grammarlens}[title={What you've built}]
One more word for this chapter.
\end{grammarlens}

\subsection*{Your turn}
\begin{itemize}
\raggedright
  \item \emph{Say it:} \emph{die Umwelt}
  \item \emph{Say it:} \emph{die Umwelt}, once more
  \item \emph{Say it:} say \emph{die Umwelt}
  \item \emph{Recall:} say the German for a surprise, then the German for the surroundings, then say \emph{die Umwelt} again
\end{itemize}

\begin{tcolorbox}[breakable,colback=teal!4,colframe=teal!35!black,title={Before you move on}]
What does die Umwelt mean? (\textquotedblleft{}The environment\textquotedblright{}.) Say it once more.
\end{tcolorbox}
\section[der Unterricht]{der Unterricht — lessons, teaching}

\label{lesson:GE-C241-unterricht}



\begin{quote}
Before the new one: say the German for the surroundings, then the German for the environment.
\end{quote}

\subsection*{You'll want to know: der Unterricht}
\textbf{der Unterricht} — \textquotedblleft{}lessons, teaching\textquotedblright{}.

\begin{cousinweb}
\textbf{Unterricht} is the noun of a verb built from the word for \textquotedblleft{}under\textquotedblright{} + the verb \textquotedblleft{}to direct, to set right\textquotedblright{}, Old High German \emph{rihten} (EtymDB 2.1): instruction.
\end{cousinweb}

\begin{grammarlens}[title={What you've built}]
One more word for this chapter.
\end{grammarlens}

\subsection*{Your turn}
\begin{itemize}
\raggedright
  \item \emph{Say it:} \emph{der Unterricht}
  \item \emph{Say it:} \emph{der Unterricht}, once more
  \item \emph{Say it:} say \emph{der Unterricht}
  \item \emph{Recall:} say the German for the surroundings, then the German for the environment, then say \emph{der Unterricht} again
\end{itemize}

\begin{tcolorbox}[breakable,colback=teal!4,colframe=teal!35!black,title={Before you move on}]
What does der Unterricht mean? (\textquotedblleft{}Lessons, teaching\textquotedblright{}.) Say it once more.
\end{tcolorbox}
\section[der Unterschied]{der Unterschied — a difference}

\label{lesson:GE-C241-unterschied}



\begin{quote}
Before the new one: say the German for the environment, then the German for lessons.
\end{quote}

\subsection*{You'll want to know: der Unterschied}
\textbf{der Unterschied} — \textquotedblleft{}a difference\textquotedblright{}.

\begin{cousinweb}
\textbf{Unterschied} is the word for \textquotedblleft{}under, between\textquotedblright{} + the verb \textquotedblleft{}to separate\textquotedblright{}, Old High German \emph{sceidan} (EtymDB 2.1): what sets two things apart.
\end{cousinweb}

\begin{grammarlens}[title={What you've built}]
One more word for this chapter.
\end{grammarlens}

\subsection*{Your turn}
\begin{itemize}
\raggedright
  \item \emph{Say it:} \emph{der Unterschied}
  \item \emph{Say it:} \emph{der Unterschied}, once more
  \item \emph{Say it:} say \emph{der Unterschied}
  \item \emph{Recall:} say the German for the environment, then the German for lessons, then say \emph{der Unterschied} again
\end{itemize}

\begin{tcolorbox}[breakable,colback=teal!4,colframe=teal!35!black,title={Before you move on}]
What does der Unterschied mean? (\textquotedblleft{}A difference\textquotedblright{}.) Say it once more.
\end{tcolorbox}
\section[die Ursache]{die Ursache — a cause}

\label{lesson:GE-C241-ursache}



\begin{quote}
Before the new one: say the German for lessons, then the German for a difference.
\end{quote}

\subsection*{You'll want to know: die Ursache}
\textbf{die Ursache} — \textquotedblleft{}a cause\textquotedblright{}.

\begin{cousinweb}
\textbf{Ursache} is \emph{ur-}, \textquotedblleft{}original, first\textquotedblright{}, + \emph{Sache}, \textquotedblleft{}matter\textquotedblright{} $\leftarrow$ Old High German \emph{sahha} (EtymDB 2.1): the first matter, the cause.
\end{cousinweb}

\begin{grammarlens}[title={What you've built}]
One more word for this chapter.
\end{grammarlens}

\subsection*{Your turn}
\begin{itemize}
\raggedright
  \item \emph{Say it:} \emph{die Ursache}
  \item \emph{Say it:} \emph{die Ursache}, once more
  \item \emph{Say it:} say \emph{die Ursache}
  \item \emph{Recall:} say the German for lessons, then the German for a difference, then say \emph{die Ursache} again
\end{itemize}

\begin{tcolorbox}[breakable,colback=teal!4,colframe=teal!35!black,title={Before you move on}]
What does die Ursache mean? (\textquotedblleft{}A cause\textquotedblright{}.) Say it once more.
\end{tcolorbox}
\section[das Verhalten]{das Verhalten — behaviour}

\label{lesson:GE-C241-verhalten}



\begin{quote}
Before the new one: say the German for a difference, then the German for a cause.
\end{quote}

\subsection*{You'll want to know: das Verhalten}
\textbf{das Verhalten} — \textquotedblleft{}behaviour\textquotedblright{}.

\begin{cousinweb}
\textbf{Verhalten} is \emph{ver-} + the verb \textquotedblleft{}to hold\textquotedblright{}, Old High German \emph{haltan}, English \emph{hold} (EtymDB 2.1): how you hold yourself, behaviour.
\end{cousinweb}

\begin{grammarlens}[title={What you've built}]
One more word for this chapter.
\end{grammarlens}

\subsection*{Your turn}
\begin{itemize}
\raggedright
  \item \emph{Say it:} \emph{das Verhalten}
  \item \emph{Say it:} \emph{das Verhalten}, once more
  \item \emph{Say it:} say \emph{das Verhalten}
  \item \emph{Recall:} say the German for a difference, then the German for a cause, then say \emph{das Verhalten} again
\end{itemize}

\begin{tcolorbox}[breakable,colback=teal!4,colframe=teal!35!black,title={Before you move on}]
What does das Verhalten mean? (\textquotedblleft{}Behaviour\textquotedblright{}.) Say it once more.
\end{tcolorbox}
